\documentclass[10pt,a4paper]{amsart}
\usepackage[margin=19mm]{geometry}
\usepackage{amsmath,amssymb,mathtools}
\usepackage[scr=rsfs,cal=euler]{mathalfa}
\usepackage{xfrac}
\usepackage[unicode,hidelinks,bookmarks=false]{hyperref}
\newcommand{\ie}{i.e.\ }
\newcommand{\cf}{cf.\ }
\newcommand{\resp}{respectively\ }
\newcommand{\Subsection}{\subsection}
\providecommand{\nobreakdash}{\nobreak\hbox{-}}
\setlength{\emergencystretch}{2em}
\multlinegap0pt
\usepackage{bm}
\usepackage{bbm}
\usepackage{dsfont}
\usepackage{xspace}
\usepackage{cancel}
\usepackage{mathtools}
\usepackage{amsthm}
\makeatletter
\@ifundefined{definition}{\newtheorem{definition}{Definition}}{}
\@ifundefined{lemma}{\newtheorem{lemma}[definition]{Lemma}}{}
\@ifundefined{theorem}{\newtheorem{theorem}[definition]{Theorem}}{}
\makeatother
\DeclareMathOperator{\herm}{Herm}
\def\D{{\mathbb D}}
\def\N{{\mathbb N}}
\title[Grothendieck's theorem for Bessel sequences]{Grothendieck's theorem for Bessel sequences}
\author{Lukas Liehr, Mitchell A. Taylor, Peiyang Yu}
\date{Source: arXiv:2608.12280v1 (CC BY 4.0)}
\begin{document}
\maketitle

\noindent\emph{Source and licence.} Grothendieck's theorem for Bessel sequences. Version arXiv:2608.12280v1 (CC BY 4.0), distributed under the Creative Commons Attribution 4.0 International licence, as recorded in the arXiv licence metadata for this exact version and shown on the abstract page https://arxiv.org/abs/2608.12280v1. Full rights notice: licenses/arxiv-2608.12280-CC-BY-4.0.txt.\par\smallskip

\noindent\textit{Editorial scope.} Counted unit: the Theorem whose source label is \texttt{thm:countable\_gram\_realization} in the article (source0.tex lines 536--542, source label \texttt{thm:countable\_gram\_realization}), delivered together with the article's own complete original proof of it (source0.tex lines 543--627, one \texttt{proof} environment of 85 lines). No mathematics is written, repaired or re-derived: statement and proof are the article's own text line for line. Required context retained uncounted: thm:main1:finite (lines 478--484); lem:probability\_realization (lines 525--534). Every definition, display and label of those lines is retained unchanged. Omitted from this package and not redistributed: 1--477, 485--524, 535--535, 628--896. Nothing inside a retained excerpt is omitted or altered: every retained body is delivered exactly as the article's lines are written, so no omission receipt of any kind accompanies this package. Citations: the retained excerpts carry no citation command at all, so no bibliography entry of the article is transcribed and no reference is redistributed. Reference adaptation: none was needed and none was made; every reference inside the delivered unit resolves inside it, so no unresolved reference or citation is produced. Printed slips kept literal: no printed word, symbol, hypothesis or number of the counted statement or of its proof was altered. Layout and class: the article's own class is \texttt{amsart} with packages including amsmath, amsthm, amssymb, amscd, accents, amsfonts, mathtools, url, mathrsfs, dsfont, datetime, hyperref, xcolor, cite; the wrapper is the lane's pinned \texttt{amsart} editorial wrapper at 10pt on A4, which restates the theorem-like environments the retained text needs and adds the article's own macro definitions that the retained text needs (\texttt{\textbackslash D}, \texttt{\textbackslash N}, \texttt{\textbackslash herm}), with extra wrapper packages (bm, bbm, dsfont, xspace, cancel, mathtools). The delivered numbering is the unit's own. Assets: no retained span contains a figure environment, a graphics-inclusion command, a PDF-page inclusion, a table or a TikZ picture; no image file is copied into this package, the package declares no asset dependency and no image file is redistributed with it. Licence: version arXiv:2608.12280v1 of this article is distributed under the Creative Commons Attribution 4.0 International licence, as recorded in the arXiv licence metadata for this exact version and shown on the abstract page https://arxiv.org/abs/2608.12280v1. A full rights notice is delivered at licenses/arxiv-2608.12280-CC-BY-4.0.txt. Characters: every retained excerpt is pure ASCII, so the delivered engine, XeLaTeX with its default Latin Modern text font, reports no missing character.
\begin{theorem}\label{thm:main1:finite}
    Let $G=(g_{jk})_{1\le j,k\le N} \in \herm_N$ satisfy $0\leq G\leq I_N$ in the standard order. Then there exist $\{ f_j \}_{1\leq j\leq N} \subseteq L^\infty([0,1])$ with $\| f_j \|_{L^\infty([0,1])} \leq 1$ for all $j$ and a constant $0 <K< \infty$ such that
    \begin{equation}
        \langle x_j,x_k \rangle = K \int_0^1 f_j(x) \overline{f_k(x)} \, dx, \quad j,k \in \N.
    \end{equation}
    Moreover, one can choose $K=1$ and this constant is optimal.
\end{theorem}

\begin{lemma}[Realization of a probability measure]
\label{lem:probability_realization}
Let $X$ be a nonempty compact metric space and let $\nu$ be a Borel
probability measure on $X$. Then there exists a Borel measurable map
$T:[0,1]\to X$ such that
\[
\lambda\bigl(T^{-1}(B)\bigr)=\nu(B)
\]
for every Borel set $B\subseteq X$.
\end{lemma}

\begin{theorem}\label{thm:countable_gram_realization}
    Let $H=(h_{jk})_{j,k\in \N}$ be a family of complex numbers such that, for all $N\in \N$, $H_N\coloneqq(h_{jk})_{1\leq j,k\leq N}\in \herm_N$ and $0\le H_N\le I_N$ in the standard order. Then there exists $\{ f_j \}_{j\in N} \subseteq L^\infty([0,1])$ with $\| f_j \|_{L^\infty([0,1])} \leq 1$ for all $j$ and a constant $0 <K< \infty$ such that
    \begin{equation}
        h_{jk} = K \int_0^1 f_j(x) \overline{f_k(x)} \, dx, \quad j,k \in \N.
    \end{equation}
    Moreover, one can choose $K=1$ and this constant is optimal.
\end{theorem}

\begin{proof}
Let us define $Z\coloneqq\D^{\N}$, so that $Z$ is a compact metrizable space in the product topology. For example, a compatible metric can be defined as
\begin{equation*}
d((y_n)_{n\in\N},(z_n)_{n\in\N})
\coloneqq
\sum_{n=1}^\infty
2^{-n}\min\{1,|y_{n}-z_{n}|\}.
\end{equation*}
Let $\mathcal P(Z)$ denote the set of Borel probability measures
on $Z$, equipped with the topology of weak convergence. This space
is compact. Indeed, by the Riesz--Markov--Kakutani representation theorem, $\mathcal P(Z)$ may be identified with the set of positive
functionals $L\in C(Z)^*$ satisfying $L(\mathds{1})=1$. This set is weak-*
closed in the unit ball of $C(Z)^*$, and hence is compact by the
Banach--Alaoglu theorem.

For $j,k\in J$, define
\begin{equation*}
\mathcal C_{jk}
=
\left\{
\nu\in\mathcal P(Z):
\int_X z_j\overline{z_k}\,d\nu(x)=h_{jk}
\right\}.
\end{equation*}
Since $z\longmapsto z_j\overline{z_k}$ is continuous on $Z$, $\mathcal C_{jk}$ is closed in
$\mathcal P(Z)$. We now show that the family $\{\mathcal C_{jk}:j,k\in \N\}$ has the finite intersection property. Let $\mathcal S\subseteq \N\times \N$ be finite, and let $F\subseteq \N$ be the finite set consisting of all indices occurring in the pairs in $\mathcal S$. Let us define the Gram matrix of $\{ x_j \}_{j \in F}$ as 
\begin{equation*}
    H_F= (h_{jk})_{j,k\in F}\coloneqq (\langle x_j, x_k\rangle)_{j,k\in F}.
\end{equation*}
By our assumption on $H$, $0\le H_F\le I_F$. Applying Theorem~\ref{thm:main1:finite} gives functions $\{a_j^{(F)}\}_{j\in F}\subseteq L^\infty([0,1])$ such that
\begin{equation*}
\|a_j^{(F)}\|_{L^\infty([0,1])}\le1,
\qquad
\int_0^1
a_j^{(F)}(t)\overline{a_k^{(F)}(t)}\,dt
=
h_{jk},
\qquad j,k\in F.
\end{equation*}
After changing these functions on a common null set, we may assume
that
\begin{equation*}
\forall t\in [0,1],\, \forall j\in F:\quad\left|a_j^{(F)}(t)\right|\le1.
\end{equation*}
Let $\eta_F$ be the distribution on
$\D^{\,F}$ of the random vector
\begin{equation*}
    t\longmapsto
    \bigl(a_j^{(F)}(t)\bigr)_{j\in F},
\end{equation*}
and let $\iota_F:\D^{\,F}\to Z$ be the continuous map which leaves the coordinates in $F$
unchanged and sets every coordinate outside $F$ equal to zero. Define
\begin{equation*}
\nu_F\coloneqq(\iota_F)_\#\eta_F.
\end{equation*}
Then
\begin{equation*}
    \int_X z_j\overline{z_k}\,d\nu_F(z)
    =
    \int_{\D^F} z_j\overline{z_k}\, d\eta_F(z)
    = h_{jk},
    \qquad j,k\in F.
\end{equation*}
Consequently, $\nu_F\in
\bigcap_{(j,k)\in\mathcal S}\mathcal C_{jk}$. This proves the finite intersection property.

Since $\mathcal P(Z)$ is compact and $(\mathcal C_{jk})_{j,k\in\N}$ are closed, there exists $\nu\in\bigcap_{j,k\in J}\mathcal C_{jk}$. Lemma~\ref{lem:probability_realization} then gives a Borel map
\begin{equation*}
T:[0,1]\to Z
\end{equation*}
whose distribution is $\nu$. Let $\pi_j:Z\to\D$
be the $j$-th coordinate map and define
\begin{equation*}
f_j(t)=\pi_j(T(t)).
\end{equation*}
Then $|f_j(t)|\le1$ for every $t$ and
\begin{equation*}
\int_0^1 f_j(t)\overline{f_k(t)}\,dt
=
\int_Z z_j \overline{z_k}\, d\nu(z)
=
h_{jk},
\qquad j,k\in J.\qedhere
\end{equation*}
\end{proof}

\end{document}
